
% ════════════════════════════════════════════
% 具体分析（5.3.1.具体分析.tex，三级子块）
% ════════════════════════════════════════════
\subsubsection{具体分析}

针对问题三，本问属于滚动预测下的随机调度与调整优化类问题，采用滚动时域线性规划方法，在 0:00、6:00、12:00 与 18:00 四个预报时刻依次制定计划购电与调整购电策略，并引入违约惩罚与超额电价机制。

本文先将附件三中四时刻发布的小时光伏预报插值到 10 分钟分辨率，形成可逐间隔使用的功率序列；在 0:00 依据最新预报制定全天计划购电策略；在 6:00、12:00、18:00 依据不断更新的预报对剩余时段进行滚动调整，其中计划购电量高于调整购电量的部分按交易电价 50\% 计违约费用，调整购电量高于计划购电量的部分按交易电价 1.5 倍计价；最后以实际供电缺口核算 5 倍电价的紧急购电费用，得到全年总购电费用并分析引入其他预报时刻调整的收益（1 月为预热期，费用汇总仅统计 2 月 1 日—12 月 31 日）。具体分析流程如图\ref{fig:analysis3_flow}所示。

\begin{figure}[H]
    \begin{center}
        \begin{tikzpicture}[x=1cm, y=0.75cm,
            flowhead/.style={rectangle, rounded corners=3pt, draw=orange!80!black, fill=orange!55!black, text=white, font=\bfseries, minimum width=7.2cm, minimum height=1.1cm, align=center, inner sep=2pt, line width=0.7pt},
            flowbox/.style={rectangle, rounded corners=2.5pt, draw=orange!70!black, fill=orange!14, minimum width=7.2cm, minimum height=1.0cm, align=center, inner sep=2pt, line width=0.6pt},
            flowend/.style={rectangle, rounded corners=2.5pt, draw=orange!85!black, fill=orange!45!black, text=white, font=\bfseries, minimum width=7.2cm, minimum height=1.0cm, align=center, inner sep=2pt, line width=0.7pt},
            flowarrow/.style={thick, -{Stealth}, orange!60!black}
        ]
            \node[flowhead] at (0, 0) (h) {问题三：滚动预报与调整优化};
            \node[flowbox] at (0, -1.8) (s1) {\textbf{光伏预报插值}\\\footnotesize 四时刻小时预报插值到 10 分钟};
            \node[flowbox] at (0, -3.6) (s2) {\textbf{0:00 制定计划}\\\footnotesize 依最新预报制定全天购电计划};
            \node[flowbox] at (0, -5.4) (s3) {\textbf{6:00 滚动调整}\\\footnotesize 更新预报并调整剩余时段购电};
            \node[flowbox] at (0, -7.2) (s4) {\textbf{12:00 滚动调整}\\\footnotesize 更新预报并调整剩余时段购电};
            \node[flowbox] at (0, -9.0) (s5) {\textbf{18:00 滚动调整}\\\footnotesize 更新预报并调整剩余时段购电};
            \node[flowbox] at (0, -10.8) (s6) {\textbf{分段惩罚计价}\\\footnotesize 不足按 50\% · 超额按 1.5 倍};
            \node[flowbox] at (0, -12.6) (s7) {\textbf{紧急购电核算}\\\footnotesize 实际供电缺口按 5 倍电价计费};
            \node[flowend] at (0, -14.4) (s8) {2.1—12.31费用汇总\\\footnotesize 输出总购电费用与调整收益};
            \draw[flowarrow] (h) -- (s1);
            \draw[flowarrow] (s1) -- (s2);
            \draw[flowarrow] (s2) -- (s3);
            \draw[flowarrow] (s3) -- (s4);
            \draw[flowarrow] (s4) -- (s5);
            \draw[flowarrow] (s5) -- (s6);
            \draw[flowarrow] (s6) -- (s7);
            \draw[flowarrow] (s7) -- (s8);
        \end{tikzpicture}
        \caption{问题三分析流程图}
        \label{fig:analysis3_flow}
    \end{center}
\end{figure}
